\section{Introduction}

Around the world, ecosystems are experiencing alarming declines in biodiversity.
Conservation efforts have shown encouraging results in some cases,\cite{bd_1,bd_2,bd_3}
though preservation measures
are constrained by available resources.\cite{conservation_funding}
Appropriate action must be informed by accurate data (e.g., population size and density), stressing
the urgent need for larger scale biomonitoring capabilities.\cite{conservation_data}

Through the $\mathrm{21}^{\mathrm{st}}$ century, the usage of camera traps for biomonitoring has
increased exponentially.\cite{ct_usage}
These are autonomous cameras, typically motion triggered, that are deployed in natural habitats to
capture footage of wildlife.
Camera traps have proven invaluable for conservation efforts,\cite{ct_success} providing the
infrastructure for cheap, large-scale data collection from which practitioners can form conclusions about population
health and allocate resources appropriately in order to maximise results.\cite{ct_usage}

Knowledge of population density and size (abundance) of animal populations is extremely important.
For species where individuals are identifiable by unique markings, capture/recapture style methods
can be applied in the analysis of camera trap video to derive estimates;\cite{capture_recapture}
however, this is not so straightforward for unmarked species.
In these cases, estimates necessitate animal-to-camera distance data.\cite{ctds}
Generally, camera traps are monocular, meaning they use a single lens to capture imagery and, therefore,
do not (directly) convey depth information.
This highlights a significant problem.
To obtain estimates of population density and abundance, animal-to-camera distances must be estimated.
Here, the conventional approach involves a slow, manual process by which analysts compare animal
observations in camera trap video to reference material that depicts objects at the respective camera
location where their distance to the camera is known.\cite{HAUCKE2022101536}

It follows that this constitutes a major bottleneck in conservation practices, where obtaining
estimates for population density and abundance is only achievable though an inefficient, subjective
and resource-intensive process.
An alternative method, where manual distance estimation is avoided, would allow for a better use
of resources in conservation efforts, allowing for broader biomonitoring capabilities across
ecosystems.

Herein, an automated distance estimation pipeline is applied to a real-world camera trap
dataset capturing chimpanzees in the West African rainforest.
The distance estimates obtained are then used to estimate the local chimpanzee population
density and abundance, establishing a full pipeline for the estimation of population density and
abundance from raw camera trap video.

This application of monocular depth estimation (MDE) for the estimation of population density
is a virtually unexplored area of research.
To date, there remains only a single example of a similar study in the
literature.\cite{full_pipeline_study}
The following study aims to assess the feasibility of applying a MDE pipeline to a dataset
representative of a different environment, analysing and explaining trends in distance,
density and abundance data generated using different parameterisable configurations of
the distance estimation pipeline.
It aims assess the accuracy of distance estimates and subsequent density and abundance estimates
with respect to those obtained using the manual approach, determining which configuration
of the pipeline should be used to obtain the most accurate estimates.
Finally, this study aims to evaluate limitations of this approach and identify ways in which it
may be refined.
